%! TEX root = ../main.tex
\documentclass[main]{subfiles}

\begin{document}
\section{むすび}
本稿では，DCASE2024 Task 4の評価基準モデルを基盤に，半教師あり学習手法 Confident Mean Teacherとドメイン汎化手法 MixStyleの改善の統合によって，音響イベント検出の精度向上を図った． 
結果として，負例を活用したCMTは，従来のMTや単純なCMTと比較して精度向上を達成した．
また，周波数ごとの重要度に基づくMixStyleの統合が有効であることを再確認し，軽量なCNNを用いた場合でもドメイン汎化の効果が得られることを示した．\par

今後の課題として，CMTとMixStyleを併用した際に観測された，ベストモデル到達エポック数の減少（早期収束または不安定性）に関する詳細な解析や，
他の半教師あり学習手法との性能比較，モデルの軽量化，時間推定タスクの上位レイヤ委譲などが挙げられる．\par
\end{document}